\scanpage{4}
\uline{Proposition} Soit $X$ un espace loc. noeth. dont toute partie fermée irréductible a exactement un point générique. \deleted{Soit} $F$ un faisceau sur $X$, satisfaisant la condition suivante : pour tt $x\in X$, il existe un \deleted{partie} \inserted{voisinage} ouvert \deleted{non vide} $U$ de $\overset{\deleted{\text{--}}}{x}$ tel que pour \unsure{ts|tt} $y\in\bar{x}\cap U$, l'application naturelle $F_y\to F_x$ soit \uline{injective}. \deleted{Alors}

\deleted{Soit} Ordonnons $X$ par $y\unsure{\leq|\leqq}x\Leftrightarrow\bar{y}\supset\bar{x}$ i.e. $\bar{y}\ni x$, soit $\widetilde F$ le \deleted{faisceau} système \deleted{inductif} \inserted{projectif} sur $X$ défini par $F$, considérons l'application naturelle
\[
\Gamma(X,F)\unsure{\text{---}|\longrightarrow}\varprojlim\widetilde F
\]
Cette application est \uline{bijective}.

La question est locale. On peut donc supposer $X$ noethérien. \deleted{Soit $U$ un ouvert de} \inserted{Il suffit de prouver} que $\longrightarrow$ est surjective, soit $(\xi_x)_{x\in X}\in\varprojlim F$.

\uline{Lemme 1} Soit $s$ une section de $F$ sur $U$. Alors l'ens des pts $x\in U$ tels que $s_x=\xi_x$ est \uline{ouvert}.

S'il contient $x$, il contient les $y\leq x$ i.e. \unsure{tels que|les pour} $\bar{y}\ni x$, il suffit de montrer qu'il est \uline{constructible}, donc que s'il contient $x$, il contient \deleted{\unsure{les|le}} un ouvert non vide de $\bar{x}$. Or il existe un ouvert non vide $U$ de $\bar{x}$ tel que $y\in\unsure{\mathring U|\tilde U|U}\Rightarrow \deleted{\unsure{O|\xi}}F_y\to F_x$ injectif, d'où $s_y=\xi_y$. O.K. -- La démonstration est maintenant \unsure{triviale|finie}.
% archive sheet number: 3
